We will present a proof of \autoref{thm:rga-ra-comp} by \cite{obmRaGreedy} which goes in three phases, first proving simple properties, then more 3 complex lemmas and finally using those to prove the theorem with the help of linear programming
%lots of problems here as the authors decided to define neither Greedy nor Sigma as well as going out of their way to not distinguish between U and V and defining both as integers
\newcommand*\xigma{\pi}

\subsection{basic properties}

We will begin by proving several properties.
\begin{lemma}[Prefix]
	\label{ra-greedy-prefix}
	$\forall\xigma_1,\xigma_2\in\mathrm{Perm}([n]),t\in[n]|(\forall 0\leq k\leq t| \xigma_1(k)=\xigma_2(k))\\ \implies\text{runs on }\xigma_1\ \text{and}\ \xigma_2\ \text{are identical}$
\end{lemma}
\begin{proof}
	This lemma is relatively trivial. Since the algorithm is deterministic and both permutations are identical the runs will be identical as well.
\end{proof}

Let $\Omega_p$ be a set of permutations where all vertices but $v_p$ have fixed positions, $\xigma_k\in\Omega_p$ be the permutation where $x_p$ is on position $k$, $G_s(t)$ be the set of offline vertices matched in time $1,\dots,t$ by the algorithm on $\xigma_s$

\begin{lemma}[Monotonicity]
	\label{ra-greedy-monotone}
	$\forall 1\leq s<m\leq n$:
	\begin{itemize}
		\item $\forall 1\leq t\leq s-1$ $G_m(t)=G_s(t)$
		\item $\forall s-1\leq t\leq m-1$ $G_m(t)\subseteq G_s(t+1)$
	\end{itemize}
\end{lemma}

Another way of phrasing this lemma is to say that, if a vertex is moved up in the permutation then all offline vertices matched previously by time $t$ are expected to be matched by time $t+1$

\begin{proof}
	The first part is a direct consequence of \autoref{ra-greedy-prefix}. Proof of the second part will be by induction over time $t$, starting with $s-1$ which is directly implied by the first part.\\
	Assuming the lemma holds at time $k$, $G_m(k)\subseteq G_s(k+1)$. Notice that $\xigma_m(k+1)=\xigma_s(k+2)$. Denote $v:=v_{\xigma_m(k+1)}$, matched to $u$ on a run on $\xigma_m$. Suppose $u\not\in G_s(k+2)$. Then $u$ was available when the algorithm was matching $v$ on $\xigma_s$, but $v$ was instead matched to $u'\not\in G_s(k+1)$. By induction hypothesis $u'\not\in G_m(k)$ implying $u,u'$ available at $t=k+1$ in run $\xigma_m$, the algorithm should have prioritized $u'$ over $u$, contradiction implying $u\in G_s(k+2)$ implying step condition for $t=k+1$
\end{proof}

As context for the last lemma, we need to establish the following: what it means to miss a vertex, what is a good match and what is a bad match.

A vertex is considered missed if it remains unmatched by the end of the run

\[\miss_\xigma(s,p)=\begin{cases}
	1 & \xigma(s)=p\land\ v_p\text{ unmatched at the end run on }\xigma\\
	0 & \text{ else}
\end{cases}\]

if $v_p$ was missed then $u_p$ must have been matched to some $v_{p'}$ and $v_{p'}$ was prioritized before $v_p$. We define this match (between $u_p$ and $v_{p'}$) as a bad match since in the end it caused $v_p$ to be missed

\[\bad_\xigma(s,q)=\begin{cases}
	1 & u_q\text{ is matched to }v_{\xigma(s)}\land \xigma^{-1}(q)>s\\
	0 & \text{ else}
\end{cases}\]

in contrast, a good match occurs when $v_{p'}$ is prioritized after $v_p$. In that situation we have a guarantee that $v_p$ wasn't missed since $u_p$ and $v_p$ are connected

\[\good_\xigma(s,q)=\begin{cases}
	1 & u_q\text{ is matched to }v_{\xigma(s)}\land \xigma^{-1}(q)\leq s\\
	0 & \text{ else}
\end{cases}\]

\begin{lemma}[Partitioning]
	\label{ra-greedy-partition}
	$\forall p,\Omega_p$, during a run on $\xigma_n$. If $u_p$ is matched to $v_t$ then
	\begin{align}
		\forall x\in[1,t]: & \sum\limits_{1\leq s\leq t+1} \good_{\xigma_x}(s,p)=1
		\forall x\in[t+1,n]: & \bad_{\xigma_x}(t,p)=1
	\end{align}
\end{lemma}

\begin{proof}
	1. from \autoref{ra-greedy-monotone} we know that $\forall x\in[1,t]$, in a run on $\xigma_x$ is matched to $v_k$ for some $k\in[x,t+1]$ $u_p$ is matched to $v_t$ implying inequality 1
	2. from \autoref{ra-greedy-prefix} we know that $\forall x\in[t+1,n]$, in a run on $\xigma_x$ $u_p$ is matched to some $v_k$ for $k\in[x,t+1]$ implying inequality 2
\end{proof}

\subsection{bad/good/miss properties}

The second part is to prove three helper lemmas about $\miss,\bad,\good$ with the basic properties

\begin{lemma}
	\label{ra-greedy-bgm-1}
	$\forall\xigma\in\mathrm{Perm}([n]),t\in[n]$\\
	$\sum\limits_{1\leq p\leq n}\miss_{\xigma}(t,p)+\sum\limits_{1\leq p\leq n}\bad_{\xigma}(t,p)+\sum\limits_{1\leq p\leq n}\good_{\xigma}(t,p)=1$
\end{lemma}
\begin{proof}
	this lemma is rather trivial since a vertex is either missed or not, and if it is then either $p<p',p>p'$ or $p=p'$
\end{proof}

\begin{lemma}
	\label{ra-greedy-bgm-2}
	$\forall t\in[n]$\\
	$\miss_{\xigma_t}(t,p)\leq\sum\limits_{\xigma\in\Omega_p}\sum\limits_{s<t}\frac{\bad_\xigma(s,p)}{n-s}$
\end{lemma}
\begin{proof}
	If $\miss$ is $0$ then the inequality is trivially true since $\bad$ and $n-s$ are non-negative.\\ Assume $\miss=1$ from definition, in run on $\xigma_t$ $v_p$ that is on position $t$ remains unmatched, which is only possible if $u_p$ is matched to $v_{p'}$ on position $t'$, $t'<t$ using \autoref{ra-greedy-partition} we have
	\begin{align*}
		& \sum\limits_{\xigma\in\Omega_p}\sum\limits_{s<t}\frac{\bad_\xigma(s,p)}{n-s} & \\
		\geq & \sum\limits_{\xigma\in\Omega_p}\frac{\bad_\xigma(t',p)}{n-t'}&t'<t\\
		\geq & \sum\limits_{t'+1\leq s\leq n}\frac{\bad_{\xigma_s}(t',p)}{n-t'}\\
		= & 1 & \text{Lemma }\ref{ra-greedy-partition}.2\\
		= & \miss_{\xigma_t}(t,p)
	\end{align*}
	concluding the proof
\end{proof}

\begin{lemma}
	\label{ra-greedy-bgm-3}
	$\forall t\in[n-1]$:
	$\sum\limits_{\xigma\in\Omega_p}\sum\limits_{s\leq t}\bad_\xigma(s,p)\frac{s}{n-s}\leq \sum\limits_{\xigma\in\Omega_p}\sum\limits_{s\leq t+1}\good_\xigma(s,p)$
\end{lemma}
\begin{proof}
	\begin{description}
		\item[Case 1] $u_p$ remains unmatched on a run on $\xigma_n$\\
		Let $\xigma_x\in \Omega_p$, if $u_p$ is matched in a run on $\xigma_x$ then by \autoref{ra-greedy-prefix} it has to be matched to $v_{x'},x'\leq x$ implying $\forall t | \bad_{\xigma_x}(t,p)=0$ and the lemma proof becomes trivial by non-negativity
		\item[Case 2] $u_p$ is matched in a run on $\xigma_n$ to $v_{t'}$. We divide this case into two following cases
		\begin{description}
			\item[Case 2.1] $t<t'$\\
				From \autoref{ra-greedy-partition} and \autoref{ra-greedy-bgm-1} we know which variables are $1$, and that for every $\xigma\in\Omega_p,s<t', \bad_\xigma(s,p)=0$, implying LHS is $0$ and the Lemma is again trivial
			\item[Case 2.2] finally, if $t\geq t'$
			\[\sum\limits_{\xigma\in\Omega_p}\sum\limits_{s\leq t}\bad_\xigma(s,p)\frac{s}{n-s}\]
			\[=\sum\limits_{\xigma\in\Omega_p}\bad_\xigma(t',p)\frac{t'}{n-t'}\]
			\[=t'\]
			\[=\sum\limits_{\xigma\Omega_p}\sum\limits_{s\leq t'+1}\good_\xigma(s,p)\]
			\[=\sum\limits_{\xigma\Omega_p}\sum\limits_{s\leq t+1}\good_\xigma(s,p)\]
		\end{description}
	\end{description}
\end{proof}

\subsection{final proof}

From here we have the tools to prove the \autoref{thm:rga-ra-comp}\\
Notice that the total size of the matching is $\sum_s E\left[\sum\limits_{p}\good_\xigma(s,p)\right]+\E\left[\sum\limits_{p}\bad_\xigma(s,p)\right]$. Informally, the three lemmas suggest that if there are many misses, then there must be many good and bad matches, bounding how small the expected matching can be.\\
We will start by introducing generalized miss/bad/good variables

\begin{align*}
	\miss(s)&=\underset{\xigma\sim\mathrm{Perm}([n])}{\E}\left[\sum\limits_p\miss_\xigma(s,p)\right]\\
	\bad(s)&=\underset{\xigma\sim\mathrm{Perm}([n])}{\E}\left[\sum\limits_p\bad_\xigma(s,p)\right]\\
	\good(s)&=\underset{\xigma\sim\mathrm{Perm}([n])}{\E}\left[\sum\limits_p\good_\xigma(s,p)\right]
\end{align*}

By using \autoref{ra-greedy-bgm-2} and summing over all partitions $\Omega_p$ and all $p$ we receive thee following

\[\forall t:\sum\limits_p\sum\limits_{\xigma\in\Omega}\miss_\xigma(t,p)\leq \sum\limits_p\sum\limits_{\xigma\in\Omega}\sum\limits_{s<t}\frac{\bad_\xigma(s,p)}{n-s}\]

change summation order, divide by $|\Omega|$

\[\forall t: \miss(t)\leq\sum\limits_{s<t}\frac{\bad(s)}{n-s}\]

by doing the same on \autoref{ra-greedy-bgm-3} instead we get
\[\forall t: \sum\limits_{s\leq t}\bad(s)\frac{s}{n-s}\leq\sum\limits_{s\leq t+1}\good(s)\]

with these inequalities we can rewrite the lower bound on expected matching size as a solution of a linear program

\begin{align*}
& \text{minimize}   &                 & \sum\limits_{t}\left[\bad(t)+\good(t)\right]                                 &\\
& \text{subject to} & \forall t<n     & \sum\limits_{s\leq t+1}\good(s) -\sum\limits_{s\leq t+1}\bad(s)\frac{s}{n-s} & \geq 0\\
&                   & \forall t\leq n & \good(t)+\bad(t)+\sum\limits_{s<t}\frac{\bad(s)}{n-s}                        &\geq 1\\
&                   &                 & \good(t),\bad(t),\miss(t)                                                    &\geq 0
\end{align*}

which can in turn be rewritten with $m_t=\good(t)+\bad(t),v_t=\frac{\bad(t)}{n-t}$ as the following linear program $\mathrm{LP}$

\begin{align*}
\text{minimize} & & \sum\limits_{t}m_t&\\
\text{subject to}\ & \forall t<n & \sum\limits_{s\leq t+1}m_s -n\sum\limits_{s\leq t}v_s & \geq 0\\
& \forall t\leq n & m_t+\sum\limits_{s<t}v_s&\geq 1\\
& & m_t,v_t&\geq 0
\end{align*}

$\mathrm{LP}'$ is obtained by dropping $m_{t+1}$ from equation 1\\
According to \cite{obmRaGreedy}, Appendix A,
\begin{quote}
	It can be shown that dropping the term $m_{t+1}$ from we need to first define the following: the first set of equations does not affect the LP solution by much
\end{quote}
which is supported by the fact that \cite{karp} proved a very similar statement. Therefore $m_{t+1}$ can be dropped without affecting the competitiveness bound by more than $o(1)$

The following is $\mathrm{LP}'$ and it's dual

\noindent\makebox[\textwidth][c]{%
	\begin{minipage}{\dimexpr\textwidth+4cm\relax}
    	\begin{align*}
			\text{minimize} & & \sum\limits_{t}m_t& && \text{maximize} & & \sum\limits_i\alpha_i\\
			\text{subject to}\ & \forall t<n & \sum\limits_{s\leq t}m_s -n\sum\limits_{s\leq t}v_s & \geq 0 && \text{subject to} & \forall t<n & \sum\limits_{s> t}\alpha_s -n\sum\limits_{s\geq t}\beta_s & \leq 0\\
			& \forall t\leq n & m_t+\sum\limits_{s<t}v_s&\geq 1 && & \forall t\leq n\ & \alpha_t+\sum\limits_{s\geq t}\beta_s&\leq 1\\
			& & m_t,v_t&\geq 0 && & & \alpha_t,\beta_t&\geq 0
		\end{align*}
	\end{minipage}
}

Notice that both primal and dual constraints are satisfied when all inequalities are tight. This means that the following is a feasible solution\\

\[m_t=\alpha_{n+1-t}=(1-\frac{1}{n})^{i-1}\\v_t=\beta_{n+1-t}=\frac{1}{n}(1-\frac{1}{n})^{i-1}\]

And in the above solution $\sum m_t=\sum \alpha_t$, then this solution is optimal by strong duality.

with the objective being $\sum m_t\approx n(1-\frac{1}{e})$, concluding the proof.

This result is unsurprising given the previous chapter. Indeed, if this value was any lower it would contradict the first section as, via an argument symmetric to the second section, the optimal competitiveness ratio would be lower than $1-\frac{1}{e}$, contradicting the proof from the first section. On the other hand, if this value was higher, then it would constitute a guarantee on a competitiveness ratio higher than $1-\frac{1}{e}$, contradicting the result from the second section where, on $I$, deterministic algorithms have been shown to have this exact competitiveness ratio.